\section{Introduction}
\label{sec:intro}

AI agents that use tools are rapidly becoming the dominant pattern for deploying language
models in production. A single user request may cause an agent to invoke a dozen tools:
reading email, querying a database, calling an external API, posting to a messaging system.
Each tool call crosses a trust boundary. The tool may be provided by a third party, run on
a different server, or handle regulated data. Yet the agent's orchestration layer has no
standard mechanism for verifying what policies governed the call, who authorized the
invocation, or whether the returned data places compliance obligations on everything the
agent does next.

Model Context Protocol (MCP)~\citep{mcp2024} formalizes this interaction as a JSON-RPC
protocol with standardized tool discovery and invocation semantics. MCP has seen rapid
adoption: within a year of publication, hundreds of MCP servers expose tools ranging from
filesystem access to medical record retrieval. This adoption makes the governance gap concrete.
When an agent calls \texttt{ehr.get\_patient} through an MCP server and the response contains
protected health information, nothing in the MCP protocol tells the agent, its orchestrator,
or a downstream auditor: which policy authorized this call; whether the response sensitivity
has been tracked; whether the tool definition matches what was reviewed at deployment time.

cMCP addresses this gap. It is a TEE-based gateway that sits between an MCP client and one or
more MCP servers. Every tool call passes through cMCP, which evaluates Cedar~\citep{cedar2023}
policies before dispatch and, after the call completes, appends a signed entry to a
hash-chained audit log. At session close, cMCP produces a TRACE Claim~\citep{trace2025}: a
cryptographically signed attestation that binds the session's call record to the TEE instance
that processed it and to the policy bundle in effect at startup.

The central thesis of this paper is that hardware attestation, applied at the MCP boundary,
can provide verifiable governance properties that cannot be achieved through software-only
interception. We validate six such properties through reproducible experiments, each backed by
unit tests that run in CI without hardware TEE; the experiments establish the
software-checkable behavior of each property, while the binding to physical hardware is
architectural and is stated as a limitation (Section~\ref{sec:eval}).

\subsection*{Contributions}

\begin{itemize}[leftmargin=1.5em]
  \item \textbf{cMCP architecture.} A TEE gateway design that is transparent to existing MCP
    clients and servers, requiring no protocol changes.

  \item \textbf{TEE-measured policy enforcement (Property~\ref{prop:policy-hash}).} Policy
    bundles are measured into the TEE at startup. Any runtime modification raises
    \texttt{PolicyHashMismatch} and invalidates the TRACE Claim signature.

  \item \textbf{Session sensitivity tracking (Property~\ref{prop:session-sensitivity}).}
    A monotonically increasing session sensitivity state machine detects PHI cross-boundary
    violations that per-call policy evaluation cannot see, at a quantified false-discovery cost.

  \item \textbf{Rug-pull detection (Property~\ref{prop:rug-pull}).} Tool catalog entries are
    hash-pinned at TEE startup. Any post-approval modification to a tool's description or
    input schema triggers fail-closed rejection before the call is dispatched.

  \item \textbf{TRACE Claim key binding (Property~\ref{prop:nonce}).} Each TRACE Claim
    carries a nonce formed from the RFC~7638 JWK Thumbprint of the TEE public key and a
    random salt, committed into the hardware report; the per-session report replaces the
    salt with a hash of the session's audit-chain root. A verifier re-derives the
    thumbprint from the claim's confirmation key to bind the report to the signing key, and
    the session identifier is covered by the claim signature. Detecting replay and omission
    additionally needs a freshness check and an expected session boundary.

  \item \textbf{Temporal adjacency provenance (Property~\ref{prop:temporal}).} cMCP
    records a call graph where edges represent temporal adjacency. This is a conservative
    provenance approximation: zero false negatives by construction, with a quantified false
    positive rate.

  \item \textbf{Cross-organizational attestation chains (Property~\ref{prop:cross-org}).}
    A dual-TEE protocol links independent Phase~1 (gateway) and Phase~2 (server) TRACE
    Claims through a shared session identifier and independently verifiable nonces.
\end{itemize}

\subsection*{Paper structure}

Section~\ref{sec:background} covers MCP mechanics, TEE attestation, the TRACE standard,
Cedar policies, and the threat model. Section~\ref{sec:design} describes the cMCP
architecture and its key mechanisms. Section~\ref{sec:claims} formally states the six
security properties. Section~\ref{sec:eval} presents experimental validation of each
property with quantitative results. Section~\ref{sec:related} surveys related work.
Section~\ref{sec:conclusion} concludes.
